\documentclass[11pt]{article}
\usepackage[margin=1in]{geometry}
\usepackage{enumitem}
% =============================================================================
% Preambles/header.tex — shared LaTeX/Beamer preamble for this project
%
% Use from a Beamer lecture by prepending:
%     \input{header}
% and compiling with TEXINPUTS=../Preambles:$TEXINPUTS (the /compile-latex
% skill does this automatically).
%
% PALETTE CONTRACT. Color names defined here MUST match the names in
% Quarto/theme-template.scss so Beamer and Quarto renderings use the same
% palette. The scripts/check-palette-sync.sh script enforces this.
%
% To customize: edit the HEX values below AND the matching $variable in
% Quarto/theme-template.scss, then run scripts/check-palette-sync.sh.
% =============================================================================

% -----------------------------------------------------------------------------
% Engine detection + encoding
%
% Our /compile-latex skill uses XeLaTeX, where inputenc is unnecessary and
% can produce encoding warnings. Only load inputenc under pdfTeX.
% -----------------------------------------------------------------------------
\usepackage{iftex}
\ifPDFTeX
  \usepackage[utf8]{inputenc}
\fi

\usepackage{amsmath, amssymb, amsthm}
\usepackage{booktabs}
\usepackage{graphicx}

% -----------------------------------------------------------------------------
% PALETTE — mirrors Quarto/theme-template.scss variable names.
% Load xcolor and define colors BEFORE anything that references them
% (notably hyperref, hypersetup, and the Beamer color assignments).
% -----------------------------------------------------------------------------
\usepackage{xcolor}

\definecolor{primary-blue}{HTML}{012169}      % headings, accents
\definecolor{primary-gold}{HTML}{B9975B}      % emphasis, borders
\definecolor{highlight-yellow}{HTML}{F2A900}  % markers, alerts
\definecolor{light-bg}{HTML}{E8EDF5}          % title / section backgrounds
\definecolor{jet}{HTML}{1A1A1A}               % body text

% Semantic colors (used in diagrams and annotations)
\definecolor{positive}{HTML}{15803D}          % good / observed / identified
\definecolor{negative}{HTML}{B91C1C}          % bad / problematic / bias
\definecolor{neutral}{HTML}{525252}           % reference / context

% Highlight accents (match .hi-* classes in the SCSS)
\definecolor{hi-slate}{HTML}{314F4F}
\definecolor{hi-green}{HTML}{2E8B57}
\definecolor{hi-red}{HTML}{C41E3A}

% Semantic aliases so skill templates and snippets can write
% \textcolor{accent}{...} without hard-coding "primary-blue".
\colorlet{accent}{primary-blue}
\colorlet{accent2}{primary-gold}

% -----------------------------------------------------------------------------
% Hyperref — loaded AFTER xcolor + palette so link colors resolve.
% -----------------------------------------------------------------------------
\usepackage{hyperref}
\hypersetup{colorlinks=true, linkcolor=primary-blue, urlcolor=primary-blue,
            citecolor=primary-blue, pdfborder={0 0 0}}

% -----------------------------------------------------------------------------
% Course code — set once per deck (after \input{header}, before \begin{document})
% via \coursecode{CS401}. Shown in the footer on every slide so a deck's course
% is identifiable without opening the title page. Empty by default.
% -----------------------------------------------------------------------------
\makeatletter
\newcommand{\coursecode}[1]{\gdef\@coursecodetext{#1}}
\gdef\@coursecodetext{}
\makeatother

% -----------------------------------------------------------------------------
% Beamer theme assignments (only apply under Beamer).
%
% We detect Beamer via \@ifclassloaded{beamer}, which is the documented,
% non-internal way. This must be wrapped in \makeatletter/\makeatother
% because \@ifclassloaded contains an @-catcode character.
% -----------------------------------------------------------------------------
\makeatletter
\@ifclassloaded{beamer}{%
  \usefonttheme{professionalfonts}
  \setbeamercolor{structure}{fg=primary-blue}
  \setbeamercolor{frametitle}{fg=primary-blue}
  \setbeamercolor{title}{fg=primary-blue}
  \setbeamercolor{subtitle}{fg=primary-gold}
  \setbeamercolor{author}{fg=primary-blue}
  \setbeamercolor{institute}{fg=neutral}
  \setbeamercolor{date}{fg=primary-gold}
  \setbeamercolor{itemize item}{fg=primary-blue}
  \setbeamercolor{itemize subitem}{fg=primary-gold}
  \setbeamercolor{alerted text}{fg=primary-gold}
  \setbeamercolor{block title}{fg=white, bg=primary-blue}
  \setbeamercolor{block body}{bg=light-bg}
  % Semantic block variants: block=definition/concept (blue), exampleblock=
  % worked example (green/positive), alertblock=key takeaway or caution (gold).
  % Keep to <=2 colored boxes per slide across ALL three variants (INV-7).
  \setbeamercolor{block title example}{fg=white, bg=positive}
  \setbeamercolor{block body example}{bg=light-bg}
  \setbeamercolor{block title alerted}{fg=white, bg=primary-gold}
  \setbeamercolor{block body alerted}{bg=light-bg}

  % Minimal footer: course code (left) + page X / N (right), no navigation clutter
  \setbeamertemplate{navigation symbols}{}
  \setbeamertemplate{footline}{%
    \color{neutral}\footnotesize\hspace{0.7em}\@coursecodetext\hfill
    \insertframenumber/\inserttotalframenumber\hspace{0.7em}\vspace{0.3em}}
}{}
\makeatother

% -----------------------------------------------------------------------------
% TikZ setup — libraries the snippet gallery uses
% -----------------------------------------------------------------------------
\usepackage{tikz}
\usetikzlibrary{arrows.meta, positioning, calc, decorations.pathreplacing,
                fit, shapes.geometric, backgrounds}

% Shared TikZ styles — snippets and hand-written diagrams can pick these up
% via `\begin{tikzpicture}[every node/.style={...}]` or by naming specific
% styles (e.g., `\node[dag-node] (X) at (0,0) {X};`).
\tikzset{
  % Node styles for causal DAGs and similar
  dag-node/.style = {
    draw=primary-blue, fill=light-bg, rounded corners=2pt,
    minimum width=1.2cm, minimum height=0.9cm, align=center, font=\small
  },
  decision-node/.style = {
    draw=primary-gold, fill=white, diamond, aspect=1.8,
    minimum width=1.8cm, minimum height=1.2cm, align=center, font=\small,
    inner sep=1pt
  },
  % Edge styles — observed vs counterfactual
  observed-edge/.style = {-{Latex[length=2.5mm]}, thick, draw=jet},
  counterfactual-edge/.style = {-{Latex[length=2.5mm]}, thick, draw=neutral, dashed},
  confound-edge/.style = {-{Latex[length=2.5mm]}, thick, draw=negative, dashed},
  % Data-point styles for plots
  observed-dot/.style = {fill=primary-blue, draw=primary-blue, circle, inner sep=1.2pt},
  counterfactual-dot/.style = {fill=white, draw=neutral, circle, inner sep=1.2pt}
}

% -----------------------------------------------------------------------------
% Convenience macros
% -----------------------------------------------------------------------------
\newcommand{\muted}[1]{\textcolor{neutral}{#1}}
\newcommand{\key}[1]{\textcolor{primary-gold}{\textbf{#1}}}
\newcommand{\good}[1]{\textcolor{positive}{#1}}
\newcommand{\bad}[1]{\textcolor{negative}{#1}}

% Standout frame macro: full-bleed dark background for section transitions.
% Beamer-only (uses \begin{frame}/\setbeamercolor/\usebeamerfont) -- do not
% call from an article-class file (e.g. Notes/<CODE>/*-notes.tex). \muted,
% \key, \good, \bad, and \coursecode above are plain xcolor/gdef and are
% safe to use from either class; verified when Notes/article-class support
% was added.
% Expand: \transitionslide{Your transition title}
\newcommand{\transitionslide}[1]{%
  \begingroup
  \setbeamercolor{background canvas}{bg=primary-blue}
  \begin{frame}[plain]
    \centering
    \vfill
    {\usebeamerfont{title}\color{white}\Large #1\par}
    \vfill
  \end{frame}
  \endgroup
}

% Section divider frame (Beamer-only), an alternative to \transitionslide with a
% darker two-line style: near-black (jet) background, a small white label line
% above a larger gold title, and a thin gold rule along the bottom edge.
% Expand: \sectiondivider{Section 2}{Pointers}
\newcommand{\sectiondivider}[2]{%
  \begingroup
  \setbeamercolor{background canvas}{bg=jet}
  \begin{frame}[plain]
    \vspace*{3.0cm}
    {\color{white}\ttfamily\large #1\par}
    \vspace{0.5cm}
    {\color{primary-gold}\LARGE #2\par}
    \begin{tikzpicture}[remember picture, overlay]
      \draw[primary-gold, line width=2.5pt]
        ($(current page.south west)+(0,0.1)$) -- ($(current page.south east)+(0,0.1)$);
    \end{tikzpicture}
  \end{frame}
  \endgroup
}

% -----------------------------------------------------------------------------
% Channel watermark — "Concepts That Click" (YouTube).
%
% Article-class files (CompetitiveExam Q/A sets, Notes, Assignments) get a
% light diagonal draft watermark on every page plus a centered footer naming
% the channel. Beamer decks get a subtle TikZ background watermark instead
% (the footline already carries the course code). Centralized here so every
% MCQ PDF is branded without per-file edits.
% -----------------------------------------------------------------------------
\makeatletter
\@ifclassloaded{article}{%
  \usepackage{draftwatermark}
  \SetWatermarkText{Concepts That Click}
  \SetWatermarkScale{1}
  \SetWatermarkFontSize{2cm}
  \SetWatermarkLightness{0.86}
  \SetWatermarkAngle{45}
  \usepackage{fancyhdr}
  \pagestyle{fancy}
  \fancyhf{}
  \fancyfoot[C]{\footnotesize\textcolor{neutral}{Concepts That Click~|~YouTube: \href{https://www.youtube.com/@concepts.thatclick}{@concepts.thatclick}}}
  \renewcommand{\headrulewidth}{0pt}
  \renewcommand{\footrulewidth}{0pt}
  \fancypagestyle{plain}{%
    \fancyhf{}%
    \fancyfoot[C]{\footnotesize\textcolor{neutral}{Concepts That Click~|~YouTube: \href{https://www.youtube.com/@concepts.thatclick}{@concepts.thatclick}}}%
    \renewcommand{\headrulewidth}{0pt}%
    \renewcommand{\footrulewidth}{0pt}%
  }
}{}
\@ifclassloaded{beamer}{%
  \setbeamertemplate{background}{%
    \begin{tikzpicture}[remember picture, overlay]
      \node[rotate=45, opacity=0.055, align=center]
        at (current page.center)
        {\Huge\textcolor{neutral}{Concepts That Click}};
    \end{tikzpicture}%
  }
}{}
\makeatother 
\coursecode{JKSSB}
\renewcommand{\thesection}{54.\arabic{section}}
\newtheorem{example}{Example}[section]
\newenvironment{definitionbox}[1]{%
  \par\noindent\textbf{Definition (#1).}\ \itshape
}{\par\vspace{0.3em}}
\title{Office Shortcut and Product Trap Lab --- Detailed Lecture Notes}
\author{JKSSB Computer Awareness}
\date{\today}

\begin{document}
\maketitle

\section{Why this lesson is a trap drill}
The computer section of JKSSB papers does not reward vague recognition. It
tests the \textbf{exact key} and the \textbf{exact product category}. A wrong
option is usually not random; it is the nearest neighbour to the right answer
--- the key that differs by one letter, the same function key in a different
program, or a familiar software name placed in the wrong category. The safe
method is to read every option as a claim to be tested and to keep the exact
English term: \texttt{F5}, \texttt{F7}, \texttt{Ctrl+H}, \texttt{\$A\$1}.
This lesson gathers the Word, PowerPoint, Excel, and product-category traps
that JKSSB reuses.

\section{MS Word: exact shortcut keys}
Word has a small set of shortcuts that JKSSB asks again and again. Learn each
key with its full action.

\begin{center}
\begin{tabular}{p{0.26\textwidth}p{0.62\textwidth}}
\toprule
\textbf{Shortcut} & \textbf{Action in MS Word} \\
\midrule
F7 & Spelling \& Grammar (spell check) \\
Shift+F7 & Thesaurus \\
Ctrl+H & Find and Replace \\
Ctrl+F & Find only (opens the Navigation pane) \\
Ctrl+K & Insert Hyperlink \\
Ctrl+X & Cut the selection \\
Ctrl+C & Copy the selection \\
Ctrl+V & Paste the Clipboard contents \\
\bottomrule
\end{tabular}
\end{center}

\begin{definitionbox}{Find and Replace}
A Word tool, opened with \textbf{Ctrl+H}, that searches for a word or phrase
and replaces it with another across the document.
\end{definitionbox}

The one distinction to fix first is \textbf{Ctrl+H} versus \textbf{Ctrl+F}.
\texttt{Ctrl+H} opens Find and Replace, so it can change text;
\texttt{Ctrl+F} only finds text. The second is \textbf{F7} versus
\textbf{Shift+F7}: \texttt{F7} checks spelling and grammar, while
\texttt{Shift+F7} opens the Thesaurus. A third is \textbf{Ctrl+K}, which
inserts a hyperlink, not a comment or a caption.

\section{Word views and review tools}
Word offers several views of the same document. \textbf{Print Layout} shows
the document exactly as it will appear on the printed page, with margins,
headers, and footers in place. The other common views are Web Layout, Outline,
and Draft. For exam purposes, the key phrase ``as it will appear when printed''
points to Print Layout.

\textbf{Track Changes} is Word's review tool. When it is switched on, every
insertion, deletion, and formatting change is recorded rather than applied
silently. Each marked change can later be \textbf{accepted} or
\textbf{rejected} by a reviewer. Track Changes also stores the author's name
and the time of the edit as document metadata. The important trap is that
Track Changes does \emph{not} delete the original text permanently: the
original remains available until a reviewer accepts or rejects the change
(PYQ-12).

\begin{example}
A candidate reads the option ``Track Changes permanently deletes the original
text, so it can never be restored.'' This is false. Track Changes preserves
the original text and marks the edit; the reviewer decides whether to accept
or reject it. The trap works because ``records changes'' sounds close to
``deletes the old version.''
\end{example}

\section{Word: Section Break versus Page Break}
A \textbf{Page Break} starts a new page but keeps the document in the same
section. A \textbf{Section Break} starts a new section. The distinction
matters because only a Section Break lets the pages that follow use a
\emph{different} header or footer. This is a repeated trap: the wrong option
claims that a Page Break changes the header or footer, which it does not
(PYQ-13; TRAP-13).

\begin{definitionbox}{Section Break}
A break that divides a document into sections, so each section can have its
own page setup, including a different header and footer.
\end{definitionbox}

\section{Word feature trap: Mail Merge}
\textbf{Mail Merge} combines a \textbf{main document} (the common letter or
label) with a \textbf{data source} such as an address database, and produces
many personalised copies. It is a Word feature. The trap is to confuse it
with Track Changes or AutoCorrect. The reliable clue in the question is the
pair ``combines letters with an address database'', which points to Mail
Merge and not to any other Word tool (PYQ-37; TRAP-29).

\section{PowerPoint: the show keys}
PowerPoint's presentation shortcuts are a favourite near-neighbour set.

\begin{center}
\begin{tabular}{p{0.26\textwidth}p{0.60\textwidth}}
\toprule
\textbf{Shortcut} & \textbf{Action in PowerPoint} \\
\midrule
F5 & Start the slide show from the beginning \\
Shift+F5 & Start the slide show from the current slide \\
Alt+F5 & Start the presentation in Presenter View \\
Esc & End the running slide show \\
\bottomrule
\end{tabular}
\end{center}

The classic trap is to call \texttt{F5} the spell-check key. In PowerPoint,
\texttt{F5} starts the show from the first slide; \texttt{F7} checks
spelling. \texttt{Shift+F5} begins from the slide you are on, and
\texttt{Alt+F5} opens Presenter View, which shows the speaker notes on the
presenter's screen while the audience sees only the slides. \texttt{Esc} ends
the show but does not jump to the first slide (PYQ-40; TRAP-32).

\section{Handouts, Slide Master, and orientation}
\textbf{Handouts} are a print layout that places several slides on one page
for the audience, usually with lined space for notes. The \textbf{Handout
Master} controls the layout of handouts and notes pages.

The \textbf{Slide Master} is the design template behind the whole
presentation. It controls the fonts, colours, background, and placeholder
positions of every slide that uses it. Editing the Slide Master changes all
dependent slides at once, which is exactly why it is used for a consistent
look.

\textbf{Orientation} is the direction of the page. \textbf{Portrait} is taller
than it is wide; \textbf{Landscape} is wider than it is tall. Slide
orientation is set in PowerPoint, and notes and handouts have their own
orientation setting.

\section{PowerPoint: converting to SmartArt}
\textbf{SmartArt} turns text into a visual diagram such as a process, cycle,
or hierarchy. A \textbf{bulleted list}, a \textbf{block of text}, and
\textbf{indented text} can all be converted directly to SmartArt. A
\textbf{table of data} cannot be converted directly, because a table is a
structured object rather than a plain text string; the Convert to SmartArt
command is disabled for it. To use table content, the text must first be
copied into a text box (PYQ-17).

\section{Excel: the formula palette}
Every Excel formula begins with the \textbf{=} sign (PYQ-38; TRAP-30). The
\textbf{formula palette} --- the Insert Function dialog --- is the helper that
lets you build a formula. It lists the function and shows each
\textbf{argument} in its own box, so you can fill in the values one at a
time. It is opened from the \textbf{fx} button next to the formula bar. The
palette guides formula entry; it does not resize cells or change data.

\begin{definitionbox}{Formula palette}
The Excel dialog box that lets a user select a function and enter its
arguments while building a formula.
\end{definitionbox}

\section{Excel: cell references and the \$ sign}
A cell reference tells Excel which cell to use. The reference type decides
what happens when the formula is copied to another cell.

\begin{center}
\begin{tabular}{p{0.20\textwidth}p{0.28\textwidth}p{0.40\textwidth}}
\toprule
\textbf{Reference} & \textbf{Type} & \textbf{Behaviour when copied} \\
\midrule
A1 & Relative & Both column and row shift \\
\$A\$1 & Absolute & Neither column nor row changes \\
\$A1 & Mixed & Column fixed, row shifts \\
A\$1 & Mixed & Row fixed, column shifts \\
\bottomrule
\end{tabular}
\end{center}

The \textbf{\$} sign locks the part immediately after it. In \texttt{\$A\$1}
both the column and the row are locked, so the reference is
\textbf{absolute}: it always points to the same cell no matter where the
formula is copied. \texttt{\$A1} locks the column only, and \texttt{A\$1}
locks the row only; both are \textbf{mixed} references. While editing a
reference, pressing \textbf{F4} cycles through relative, absolute, and mixed
forms.

\begin{example}
A formula in \texttt{D4} reads \texttt{=B4*C4}. Copied down to \texttt{D5}, it
becomes \texttt{=B5*C5}, because the references are relative. If instead the
formula is \texttt{=\$B\$4*\$C\$4}, it stays \texttt{=\$B\$4*\$C\$4} in
every cell, because the \texttt{\$} signs make both references absolute.
\end{example}

\section{Product-category traps: Word, Excel, Access, MS-DOS}
JKSSB often places four familiar names together and asks which one belongs to
a given category. Match each name to its exact category.

\begin{center}
\begin{tabular}{p{0.26\textwidth}p{0.58\textwidth}}
\toprule
\textbf{Program} & \textbf{Correct category} \\
\midrule
MS Word & Word processor (text documents) \\
MS Excel & Spreadsheet (cells and formulas) \\
MS Access & Database management system (tables, queries, reports) \\
MS-DOS & Disk Operating System (a command-line operating system) \\
\bottomrule
\end{tabular}
\end{center}

The trap is that all four are familiar computing names. \textbf{MS Access} is
a database, \emph{not} a spreadsheet; \textbf{MS Excel} is a spreadsheet,
\emph{not} a database; and none of Word, Excel, or Access is an operating
system. \textbf{MS-DOS} is the only operating system in the group.

\section{Exam retrieval and common confusions}
\begin{enumerate}
  \item Word \textbf{F7} is spell check; \textbf{Ctrl+H} is Find and Replace;
    \textbf{Ctrl+K} is Insert Hyperlink.
  \item \textbf{Section Break}, not Page Break, allows a different
    header/footer on later pages (TRAP-13).
  \item PowerPoint \textbf{F5} starts from the beginning, \textbf{Shift+F5}
    from the current slide, \textbf{Alt+F5} is Presenter View, and
    \textbf{Esc} ends the show (TRAP-32).
  \item A \textbf{table} cannot be converted directly to SmartArt.
  \item Excel \textbf{\$A\$1} is an absolute reference; every formula starts
    with \textbf{=}.
  \item \textbf{MS Access} is a database; \textbf{MS-DOS} is an operating
    system.
\end{enumerate}

\begin{example}
A question asks which key starts a PowerPoint slide show from the beginning.
Trace the set: \texttt{F5} begins the show, \texttt{Shift+F5} starts from the
current slide, \texttt{Alt+F5} opens Presenter View, and \texttt{Esc} ends the
show. The answer is \texttt{F5}. If the stem instead says ``from the current
slide'', the answer is \texttt{Shift+F5}.
\end{example}

\subsection*{Exercises}
\begin{enumerate}[label=\arabic*.]
  \item State the action of F7, Ctrl+H, and Ctrl+K in MS Word.
  \item Distinguish Print Layout from the other Word views.
  \item Explain what Track Changes records and why it does not permanently
    delete the original text.
  \item Distinguish a Page Break from a Section Break, and say which one
    allows a different header.
  \item Give the PowerPoint keys for starting a show from the beginning, from
    the current slide, for Presenter View, and for ending the show.
  \item Explain the difference between the relative reference A1 and the
    absolute reference \$A\$1 in Excel.
  \item Classify MS Word, MS Excel, MS Access, and MS-DOS by category.
\end{enumerate}

\subsection*{Solutions}
\begingroup\small
\begin{enumerate}[label=\arabic*.]
  \item \texttt{F7} opens Spelling \& Grammar; \texttt{Ctrl+H} opens Find and
    Replace; \texttt{Ctrl+K} inserts a hyperlink.
  \item Print Layout shows the document as it will appear on the printed
    page, with margins and headers in place. Web Layout, Outline, and Draft
    are the other common views and are not print-accurate.
  \item Track Changes records insertions, deletions, and formatting edits,
    along with author and time metadata. It marks the edits instead of
    applying them silently, so the original text is preserved until a
    reviewer accepts or rejects each change.
  \item A Page Break starts a new page within the same section; a Section
    Break starts a new section. Only a Section Break allows a different
    header or footer on the following pages.
  \item \texttt{F5} from the beginning; \texttt{Shift+F5} from the current
    slide; \texttt{Alt+F5} for Presenter View; \texttt{Esc} to end the show.
  \item A1 is relative, so both its column and row shift when the formula is
    copied. \$A\$1 is absolute, so neither the column nor the row changes; the
    reference always points to the same cell.
  \item MS Word is a word processor; MS Excel is a spreadsheet; MS Access is
    a database management system; MS-DOS is an operating system.
\end{enumerate}
\endgroup
\end{document}
